% Tarjetas visibles sólo en la edición pedagógica. Se invocan en el punto
% matemático correspondiente; el archivo no produce contenido por sí solo.

\newcommand{\HAFOProblemStateMemory}{\ifHAFOPedagogy
\begin{bloqueestudio}{Ejercicios}
\textbf{Formalismo.} Para \(A\in\mathbb R^{n\times n}\), compare
\(\dot x=Ax\) con \({}^{\mathrm C}D^q x=Ax\): solución fundamental,
composición temporal, espacio de fases y dato inicial. Verifique el caso
escalar con \(e^{\lambda t}\) y \(E_q(\lambda t^q)\).

\textbf{Interpretación.} Identifique exactamente dónde falla el semigrupo
puntual para \(q<1\).

\textbf{Defensa oral.} Explique en dos minutos por qué el estado de Caputo no
queda determinado sólo por \(x(t)\).
\end{bloqueestudio}
\fi}

\newcommand{\HAFOProblemStabilityBasins}{\ifHAFOPedagogy
\begin{bloqueestudio}{Ejercicios}
\textbf{Formalismo.} Para un Chua entero y uno fraccionario, calcule
equilibrios y espectros, aplique el criterio de estabilidad correspondiente y
defina la cuenca en el espacio de estados o de historias.

\textbf{Interpretación.} Separe estabilidad local, atracción y contacto de
cuenca.

\textbf{Defensa oral.} Responda por qué tener todos los equilibrios inestables
no demuestra que un atractor sea oculto.
\end{bloqueestudio}
\fi}

\newcommand{\HAFOProblemReturn}{\ifHAFOPedagogy
\begin{bloqueestudio}{Ejercicios}
\textbf{Formalismo.} Construya una sección para PLL o MAVPD, formule
\(P-I\) y las hipótesis de grado. Después escriba el operador de retorno
Caputo sobre historias.

\textbf{Interpretación.} Distinga punto fijo de retorno, estabilidad del ciclo
y existencia por grado.

\textbf{Defensa oral.} Explique por qué la proyección puntual de un retorno
Caputo no conserva toda la información dinámica.
\end{bloqueestudio}
\fi}

\newcommand{\HAFOProblemChaos}{\ifHAFOPedagogy
\begin{bloqueestudio}{Ejercicios}
\textbf{Formalismo.} Compare una trayectoria entera y una Caputo con la
misma matriz de observables. Separe acotación, recurrencia, atracción, caos,
cuenca y ocultedad.

\textbf{Interpretación.} Identifique qué afirmaciones son topológicas y cuáles
son sólo numéricas y finitas.

\textbf{Defensa oral.} Explique por qué una FFT o un exponente positivo aislado
no certifican simultáneamente caos, atracción y ocultedad.
\end{bloqueestudio}
\fi}

\newcommand{\HAFOProblemLocalization}{\ifHAFOPedagogy
\begin{bloqueestudio}{Ejercicios}
\textbf{Formalismo.} En un mismo sistema produzca semillas mediante dos métodos
entre DF, Machado, continuación, PP/FPP y borde. Compare residuos algebraicos,
especificación de memoria, destino y prueba de cuenca bajo el mismo presupuesto.

\textbf{Interpretación.} Determine qué aporta cada método y qué afirmación no
puede sostener por sí sola.

\textbf{Defensa oral.} Distinga con precisión semilla, trayectoria, conjunto
límite, atractor y atractor oculto.
\end{bloqueestudio}
\fi}

\newcommand{\HAFOThesisDefense}{\ifHAFOPedagogy
\begin{bloqueestudio}{Ejercicios}
\textbf{Simulacro de defensa.} Responda sin consultar el texto y después
contraste cada respuesta con las derivaciones:
\begin{enumerate}[nosep]
  \item ¿Cómo se pasa del campo a equilibrios, jacobiano, espectro y
  clasificación local?
  \item ¿Cuál es la diferencia entre atractor autoexcitado y oculto y cómo se
  ensaya su cuenca?
  \item ¿Cómo se obtiene la realización de Lur'e, \(W_q\) y el cierre de
  balance armónico?
  \item ¿Qué predicen DF y Machado y qué no certifican?
  \item ¿Por qué Caputo exige Volterra, terminal inferior e historia completa?
  \item ¿Cómo se continúa un sistema fraccionario sin reiniciar la memoria?
  \item ¿Qué significan PP, FPP--A y FPP--H y cuándo pueden fallar?
  \item ¿Qué aportan grado, Poincaré, variedades, superficies críticas y borde
  de cuenca?
  \item ¿Qué evidencia separa acotación, recurrencia, atracción, caos y
  ocultedad?
  \item Para cada caso analizado: ¿qué es previo, qué es propuesta del
  proyecto, qué se verificó y qué permanece abierto?
\end{enumerate}
\end{bloqueestudio}
\fi}

\newcommand{\HAFOCardTopology}{\ifHAFOPedagogy
\begin{bloqueestudio}{Lecturas}
\begin{itemize}[nosep]
  \item \emph{Topología y geometría diferencial con aplicaciones a la física}
  \cite{NahmadAchar2022Topologia}: coordenadas adaptadas y curvas integrales,
  §§11.3--11.5, pp.~124--125; derivadas de Lie, §§11.6--11.14,
  pp.~125--130.
\end{itemize}
\end{bloqueestudio}
\begin{bloqueestudio}{Ejercicios}
\textbf{Básicos.} Resuelva §§11.15--11.18, pp.~131--133. Calcule
\(L_fV=\nabla V\cdot f\) para \(V(x)=\|x\|^2/2\).

\textbf{Investigación.} Demuestre que \(\mathbb R\times S^1\) es el espacio
de fases del PLL y defina una distancia con
\(\theta\sim\theta+2\pi\).
\end{bloqueestudio}
\begin{bloqueestudio}{Lecturas complementarias}
Mismo libro: conexiones lineales, pp.~135--142.
\end{bloqueestudio}
\fi}

\newcommand{\HAFOCardFlow}{\ifHAFOPedagogy
\begin{bloqueestudio}{Lecturas}
\begin{itemize}[nosep]
  \item \emph{Teoría geométrica de ecuaciones diferenciales I}
  \cite{JaurezEtAl2021TGEDI}: cap.~1, p.~1; cap.~3, pp.~53--80; cap.~4,
  pp.~93--155.
  \item \emph{Teoría geométrica de ecuaciones diferenciales II}
  \cite{JaurezEtAl2021TGEDII}: cap.~8, pp.~69--98; existencia y unicidad,
  p.~73; primera variación, p.~74; dependencia del dato inicial, p.~92.
  \item \texttt{Strogatz - Nonlinear Dynamics and Chaos.pdf}
  \cite{Strogatz2015Nonlinear}: PDF~30--50 y 140--212.
\end{itemize}
\end{bloqueestudio}
\begin{bloqueestudio}{Ejercicios}
\textbf{Básicos.} Resuelva Strogatz 2.4.1--2.4.8, 2.5.1--2.5.5 y 2.6.2,
PDF~57--59. Para \(\dot x=Ax\), demuestre
\(\varphi_t(x_0)=e^{tA}x_0\) y
\(\varphi_{t+s}=\varphi_t\circ\varphi_s\).

\textbf{Investigación.} Calcule el jacobiano y el espectro del origen en un
caso entero. Explique qué falla en
Hartman--Grobman si hay valores propios sobre el eje imaginario.
\end{bloqueestudio}
\begin{bloqueestudio}{Lecturas complementarias}
\emph{Métodos topológicos en el estudio de las ecuaciones diferenciales no
lineales} \cite{Amster2021MetodosTopologicos}: apéndice A, pp.~319--327.
\end{bloqueestudio}
\fi}

\newcommand{\HAFOCardBasins}{\ifHAFOPedagogy
\begin{bloqueestudio}{Lecturas}
\begin{itemize}[nosep]
  \item \emph{Exploring Chaos} \cite{Davies2004ExploringChaos}:
  sistemas bidimensionales y fronteras de cuenca, pp.~107--148; §4.7,
  p.~134.
  \item \texttt{Alligood, Sauer \& Yorke - Chaos.pdf}
  \cite{AlligoodSauerYorke1996Chaos}: cuencas, PDF~146--165; fronteras,
  PDF~181--209.
\end{itemize}
\end{bloqueestudio}
\begin{bloqueestudio}{Ejercicios}
\textbf{Básicos.} Resuelva Alligood--Sauer--Yorke, \emph{Challenge} 4,
PDF~199--201.

\textbf{Investigación.} Para dos atractores \(A_1,A_2\), escriba
\(\mathcal B(A_i)\), su frontera común y un par de puntos para una bisección
de borde. Indique qué resultado sería sólo local y finito.
\end{bloqueestudio}
\begin{bloqueestudio}{Lecturas complementarias}
Alligood--Sauer--Yorke: \emph{Challenge} 10, PDF~444--451.
\end{bloqueestudio}
\fi}

\newcommand{\HAFOCardDegree}{\ifHAFOPedagogy
\begin{bloqueestudio}{Lecturas}
\emph{Métodos topológicos} \cite{Amster2021MetodosTopologicos}: índice de
curvas, pp.~74--82; disparo, p.~84; Brouwer y Poincaré, p.~99; grado de
Brouwer, pp.~103--145.
\end{bloqueestudio}
\begin{bloqueestudio}{Ejercicios}
\textbf{Básicos.} Resuelva los problemas del cap.~3, pp.~153--155, y calcule
\(\deg(I,B,0)\), \(\deg(-I,B,0)\).

\textbf{Investigación.} Para un mapa de retorno \(P\),
formule \(F(z)=P(z)-z\) y enumere las hipótesis necesarias para usar el grado.
\end{bloqueestudio}
\fi}

\newcommand{\HAFOCardLinearGeometry}{\ifHAFOPedagogy
\begin{bloqueestudio}{Lecturas}
\emph{Teoría geométrica de ecuaciones diferenciales II}
\cite{JaurezEtAl2021TGEDII}: primera variación, p.~74; dependencia del dato
inicial, p.~92; Grobman--Hartman, p.~94.
\end{bloqueestudio}
\begin{bloqueestudio}{Ejercicios}
\textbf{Básicos.} Resuelva Strogatz 6.5.11, PDF~206, y 6.8.6--6.8.7,
PDF~213--214. Derive \(\dot\Phi=Df(x(t))\Phi\), \(\Phi(0)=I\).

\textbf{Investigación.} Para un equilibrio de Chua, MAVPD, Kalman--Fitts o
PLL, separe subespacios estable, inestable y central.
\end{bloqueestudio}
\fi}

\newcommand{\HAFOCardPoincare}{\ifHAFOPedagogy
\begin{bloqueestudio}{Lecturas}
\begin{itemize}[nosep]
  \item \emph{Teoría geométrica de ecuaciones diferenciales II}
  \cite{JaurezEtAl2021TGEDII}: rectificación y Poincaré, pp.~113--115;
  Floquet y monodromía, pp.~121--141; Poincaré--Bendixson, p.~150.
  \item Strogatz \cite{Strogatz2015Nonlinear}: PDF~215--260 y 261--325.
\end{itemize}
\end{bloqueestudio}
\begin{bloqueestudio}{Ejercicios}
\textbf{Básicos.} Resuelva Strogatz 7.3.1--7.3.3, PDF~250--251, y demuestre
que una órbita periódica da un punto fijo del retorno.

\textbf{Investigación.} Construya una sección transversal para PLL o MAVPD y
explique un multiplicador de módulo mayor que uno.
\end{bloqueestudio}
\begin{bloqueestudio}{Lecturas complementarias}
Strogatz 8.4.3--8.4.4, PDF~313--314.
\end{bloqueestudio}
\fi}

\newcommand{\HAFOCardChaos}{\ifHAFOPedagogy
\begin{bloqueestudio}{Lecturas}
\begin{itemize}[nosep]
  \item Strogatz \cite{Strogatz2015Nonlinear}: Lorenz, PDF~326--371;
  Lyapunov y mapas, PDF~372--421; atractores extraños, PDF~446--476.
  \item Alligood--Sauer--Yorke \cite{AlligoodSauerYorke1996Chaos}:
  Lyapunov y herradura, PDF~211--247; atractores, PDF~248--289; caos en
  flujos, PDF~376--415; variedades y crisis, PDF~416--463.
\end{itemize}
\end{bloqueestudio}
\begin{bloqueestudio}{Ejercicios}
\textbf{Básicos.} Resuelva Strogatz 9.2.2--9.2.3 y 9.3.8--9.3.9,
PDF~366--368; 10.3.7--10.3.8 y 10.5.1--10.5.4, PDF~413--417.

\textbf{Investigación.} Defina por separado acotación, recurrencia, atracción,
caos y ocultedad. Construya un contraejemplo para cada implicación inválida.

\textbf{Integrador.} Repita la tabla de implicaciones para una trayectoria de
Caputo: sustituya periodicidad exacta por recurrencia proyectada y declare qué
información depende de la historia.
\end{bloqueestudio}
\begin{bloqueestudio}{Lecturas complementarias}
\begin{itemize}[nosep]
  \item \emph{Exploring Chaos} \cite{Davies2004ExploringChaos}, pp.~44--60,
  65--106 y 184--193.
  \item \emph{Chaos and Time-Series Analysis}
  \cite{Sprott2003ChaosTimeSeries}, pp.~72--183, 213--272 y 302--328.
  \item Ott \cite{Ott2002Chaos}: conjuntos caóticos no atractores y
  fronteras, pp.~168--210; crisis y transitorios, pp.~304--344.
  \item T{\'e}l--Gruiz \cite{TelGruiz2006Chaotic}: sillas, escapes,
  fronteras y caos transitorio, pp.~191--225.
  \item Devaney \cite{Devaney1989Chaotic}: PDF aprox.~29--33, 70,
  110--129.
  \item Hilborn \cite{Hilborn2000Chaos}: PDF aprox.~38--80 y 162--220.
\end{itemize}
\end{bloqueestudio}
\fi}

\newcommand{\HAFOCardFractal}{\ifHAFOPedagogy
\begin{bloqueestudio}{Lecturas complementarias}
\begin{itemize}[nosep]
  \item Falconer \cite{Falconer2014Fractal}: dimensión de caja,
  PDF~59--75; Hausdorff, PDF~76--97; sistemas dinámicos, PDF~238--264.
  \item \texttt{Barnsley - Fractals Everywhere.pdf}: métricas y topología,
  PDF~12--48; contracciones, PDF~49--129; dimensión, PDF~186--219.
  \item Mandelbrot: consulta visual, PDF~33 y 277; para pruebas use Falconer.
\end{itemize}
\end{bloqueestudio}
\begin{bloqueestudio}{Ejercicios}
\textbf{Básicos.} Resuelva un ejercicio de dimensión de caja de Falconer,
cap.~2; contraste después con el solucionario, PDF~1--21.

\textbf{Investigación.} Calcule la dimensión de una frontera de cuenca en dos
refinamientos y explique por qué el resultado no demuestra Wada.
\end{bloqueestudio}
\fi}

\newcommand{\HAFOCardFractional}{\ifHAFOPedagogy
\begin{bloqueestudio}{Lecturas}
\begin{itemize}[nosep]
  \item Oliveira \cite{Oliveira2019SolvedFractional}: Gamma, PDF~30--41;
  Mittag--Leffler, PDF~81--124; Laplace, PDF~130--137; GL, RL y Caputo,
  PDF~179--189.
  \item Herrmann \cite{Herrmann2014Fractional}: PDF~28--53 y 64--89.
  \item Das \cite{Das2020Kindergarten}: construcción y relación RL--Caputo,
  PDF~156--163; evaluación discreta GL/RL, PDF~183--185.
  \item \emph{The Analysis of Fractional Differential Equations}
  \cite{Diethelm2010AnalysisFDE}: pp.~13--73.
\end{itemize}
\end{bloqueestudio}
\begin{bloqueestudio}{Ejercicios}
\textbf{Básicos.} Oliveira: ejercicios 13--16 y 30--32,
PDF~142--172; ejercicios RL--Caputo 12, 14, 15, 20, 25, 32, 33 y 40,
PDF~191--224. Obtenga
\(\mathcal L\{t^{\beta-1}E_{\alpha,\beta}(\lambda t^\alpha)\}
=s^{\alpha-\beta}/(s^\alpha-\lambda)\).

\textbf{Investigación.} Derive la transformada de Caputo con sus datos
iniciales y compare RL, GL y Caputo en \(t^\mu\).

\textbf{Integrador.} Resuelva \(\dot x=\lambda x\) y
\({}^{\mathrm C}D^q x=\lambda x\). Compare \(e^{\lambda t}\) con
\(E_q(\lambda t^q)\), la ley de semigrupo y los criterios espectrales entero y
de Matignon. Después aproxime \({}^{\mathrm C}D^qx=-ax\) con GL/RL y resuelva
el PVI con ABM--PECE; explique por qué discretizar un operador no equivale
automáticamente a resolver el mismo PVI Caputo.
\end{bloqueestudio}
\begin{bloqueestudio}{Lecturas complementarias}
Oldham--Spanier \cite{OldhamSpanier1974Fractional}: p.~16,
pp.~47--60, 69--94 y 136. Herrmann: ejercicios 3.3, 5.1--5.2 y 6.1;
soluciones PDF~423, 428--429 y 432.
\end{bloqueestudio}
\fi}

\newcommand{\HAFOCardMemory}{\ifHAFOPedagogy
\begin{bloqueestudio}{Lecturas}
\begin{itemize}[nosep]
  \item Herrmann \cite{Herrmann2014Fractional}: no localidad y memoria,
  PDF~108--126.
  \item Das \cite{Das2020Kindergarten}: contexto analítico, PDF~41--47;
  ramas y cortes, PDF~43--44 y apéndice detallado PDF~503--507;
  inicialización, kernels y memoria desvaneciente RL, PDF~109--110 y 119--135;
  historia e inicialización RL, PDF~247--255; GL y composición, PDF~256--259;
  relación RL--Caputo, PDF~259--264; inicialización Caputo, PDF~264--273;
  seno con terminal inferior, PDF~273--274;
  Laplace con inicialización, PDF~275--291; resolución de EFD, PDF~292--310;
  EFD lineales e integrales, PDF~333--357; ecuaciones secuenciales, matrices y
  funciones alfa-exponenciales, PDF~358--371; sistemas multivariables,
  PDF~372--375.
  \item Oliveira \cite{Oliveira2019SolvedFractional}: ejercicio 26 sobre
  efecto de memoria, PDF~272--273.
  \item Diethelm \cite{Diethelm2010AnalysisFDE}: PVI Caputo de término único,
  existencia, unicidad, perturbaciones, regularidad y estabilidad,
  pp.~85--126 y 157--162.
\end{itemize}
\end{bloqueestudio}
\begin{bloqueestudio}{Ejercicios}
\textbf{Básicos.} Herrmann: ejercicio 8.1, PDF~126; solución, PDF~436.
Derive la ecuación integral de Volterra.

\textbf{Investigación.} Construya dos historias con el mismo valor terminal y
pasados distintos; muestre por qué reiniciar cambia el PVI y por qué un retorno
puntual no es, en general, un mapa de estado. Continúe una trayectoria con
historia heredada y reinicie otra desde el mismo punto; identifique el término
libre de Volterra que cambia.
\end{bloqueestudio}
\begin{bloqueestudio}{Lecturas complementarias}
Matignon \cite{Matignon1996} para el criterio angular.
Li--Chen--Podlubny \cite{LiChenPodlubny2010} para estabilidad directa y
Mittag--Leffler; no es una prueba de caos ni ocultedad.
\end{bloqueestudio}
\fi}

\newcommand{\HAFOCardExistence}{\ifHAFOPedagogy
\begin{bloqueestudio}{Lecturas}
\emph{Métodos topológicos} \cite{Amster2021MetodosTopologicos}: función de
Green, p.~229; Schauder, pp.~245--267; Leray--Schauder, pp.~273--300.
\end{bloqueestudio}
\begin{bloqueestudio}{Ejercicios}
\textbf{Básicos.} Resuelva Amster 7.9.1--7.9.2, pp.~313--316.

\textbf{Investigación.} Escriba el PVI de Caputo como
\(x=x_0+I^qf(x)\). Identifique continuidad,
compacidad y cota a priori requeridas por Schauder o Leray--Schauder; distinga
existencia de atracción y ocultedad.
\end{bloqueestudio}
\fi}

\newcommand{\HAFOCardLure}{\ifHAFOPedagogy
\begin{bloqueestudio}{Lecturas}
Slotine--Li \cite{SlotineLi1991}: realimentación y estabilidad,
PDF~119--175; función descriptiva, PDF~176--209. Para atractores ocultos:
Leonov--Kuznetsov \cite{LeonovKuznetsovVagaitsev2011,LeonovKuznetsov2013}
y localización de Chua \cite{KuznetsovEtAl2017}.
\end{bloqueestudio}
\begin{bloqueestudio}{Ejercicios}
\textbf{Básicos.} Resuelva Slotine--Li, cap.~5, PDF~207--209, y obtenga una
función descriptiva desde su integral de Fourier.

\textbf{Investigación.} Derive
\(W_q(s)=c^{\mathsf T}(s^qI-A)^{-1}b\), fije la rama de \(s^q\) y verifique
la equivalencia de \(1-W_+N=0\) y \(1+W_-N=0\).
\end{bloqueestudio}
\fi}

\newcommand{\HAFOCardMachado}{\ifHAFOPedagogy
\begin{bloqueestudio}{Lecturas}
Machado \cite{Machado2015FDF} y Barbosa--Machado
\cite{BarbosaMachado2002BacklashImpact}.
\end{bloqueestudio}
\begin{bloqueestudio}{Ejercicios}
\textbf{Básicos.} Para \(N=re^{i\theta}\), escriba todas las ramas de
\(N^\mu\) y calcule la principal en \(\mu=1/2\).

\textbf{Investigación.} Para \(N_\mu=N^\mu\), calcule por separado el
residuo auxiliar y el balance físico; determine cuándo la potencia sólo
diversifica semillas.

\textbf{Integrador.} Obtenga dos semillas de ramas distintas, intégrelas bajo
una misma especificación numérica, entera o Caputo, y compare destino y contacto con las cuencas
conocidas.
\end{bloqueestudio}
\fi}

\newcommand{\HAFOCardContinuation}{\ifHAFOPedagogy
\begin{bloqueestudio}{Lecturas}
Diethelm \cite{Diethelm2010AnalysisFDE}: ABM, pp.~195--210; esquema numérico
multi-término, pp.~211--224, con ejercicio en p.~225. Comparación de
aproximación frecuencial, predictor--corrector y Adomian en
Sun--He--Wang \cite{SunHeWang2022FractionalChaos}, pp.~27--74.
Predictor--corrector original \cite{DiethelmFordFreed2002}.
\end{bloqueestudio}
\begin{bloqueestudio}{Ejercicios}
\textbf{Básicos.} Derive los pesos predictor y corrector desde Volterra y
aplíquelos a \({}^{\mathrm C}D^qx=-x\).

\textbf{Investigación.} Escriba una homotopía \(f_\lambda\), transporte una
semilla y demuestre por qué reiniciar la memoria en cada \(\lambda\) define
otro problema. Compare dos algoritmos únicamente después de igualar operador,
terminal inferior, historia, paso y tolerancia.
\end{bloqueestudio}
\fi}

\newcommand{\HAFOCardPerpetual}{\ifHAFOPedagogy
\begin{bloqueestudio}{Lecturas}
Prasad \cite{Prasad2015Perpetual}; Dudkowski--Prasad--Kapitaniak
\cite{DudkowskiPrasadKapitaniak2015,DudkowskiPrasadKapitaniak2017,
DudkowskiPrasadKapitaniak2018};
contraejemplos \cite{NazarimehrSaediJafari2017,NagChowdhuryGhosh2020};
revisión de métodos \cite{GuanXie2025}. Das, Oliveira y Herrmann sustentan
operadores, inicialización y memoria, pero no formulan puntos perpetuos
fraccionarios: FPP--A y FPP--H se presentan como formulaciones metodológicas,
no como resultados de esas fuentes.
\end{bloqueestudio}
\begin{bloqueestudio}{Ejercicios}
\textbf{Básicos.} Para \(\dot x=f(x)\), derive
\(\ddot x=Df(x)f(x)\) y resuelva los PP de dos campos polinomiales sencillos.

\textbf{Investigación.} Resuelva \(Df(p)f(p)=0\), \(f(p)\ne0\), en uno de los
sistemas analizados. Desde Volterra obtenga
\(h^{2q}Df(p)f(p)/\Gamma(2q+1)\) sin suponer una regla de cadena no válida para
Caputo y explique
por qué su cancelación no es aceleración fraccionaria global.

\textbf{Integrador.} Para \(f(x)=Ax+d\), compare FPP--A y FPP--H. Después,
en un campo no lineal, evalúe el evento FPP--H con toda la historia y contraste
continuación heredada contra reinicio. Repita L1 con \(h,h/2,h/4\) y terminal
inferior común; un mínimo observado gráficamente no constituye un evento
validado.
\end{bloqueestudio}
\fi}

\newcommand{\HAFOCardGeometric}{\ifHAFOPedagogy
\begin{bloqueestudio}{Lecturas}
Curvas conectantes \cite{GilmoreEtAl2010Connecting,GuanXie2021Connecting};
superficies críticas \cite{GodaraEtAl2018CriticalSurfaces}; cuencas y Wada
\cite{DazaEtAl2016BasinEntropy,DazaEtAl2015Wada}.
\end{bloqueestudio}
\begin{bloqueestudio}{Ejercicios}
\textbf{Básicos.} Para un campo plano, calcule \(Df\), \(Df\,f\),
\(\det Df\) y sus conjuntos de nivel cero.

\textbf{Investigación.} Repita el cálculo en PLL o MAVPD. Formule una
bisección entre dos destinos con métrica escalada; añada un tercer destino y
describa el fallo.
\end{bloqueestudio}
\fi}

\newcommand{\HAFOCardProtocols}{\ifHAFOPedagogy
\begin{bloqueestudio}{Lecturas}
Datseris--Wagemakers \cite{datseris2022} para estimación de cuencas; entropía de
cuenca \cite{DazaEtAl2016BasinEntropy}; clasificación de cuencas
\cite{SprottXiong2015Basins}; Wada \cite{DazaEtAl2015Wada}.
\end{bloqueestudio}
\begin{bloqueestudio}{Ejercicios}
\textbf{Básicos.} Para una órbita conocida, calcule retorno, período dominante
y variación al reducir el paso.

\textbf{Investigación.} Construya una matriz
semilla\(\times\)orden--memoria\(\times\)paso\(\times\)horizonte
\(\times\)solucionador. Separe
destino, atracción, caos, contacto con equilibrios y ocultedad.
\end{bloqueestudio}
\fi}

\newcommand{\HAFOCardReliability}{\ifHAFOPedagogy
\begin{bloqueestudio}{Lecturas}
\begin{itemize}[nosep]
  \item Vallejo--Sanju{\'a}n \cite{VallejoSanjuan2017Predictability}:
  error y \emph{shadowing}, pp.~13--22; FTLE y regímenes, pp.~25--87;
  predictibilidad, pp.~91--126; algoritmo, pp.~129--136.
  \item Datseris--Parlitz \cite{DatserisParlitz2022Nonlinear}: Lyapunov y
  horizonte de predictibilidad, pp.~37--48; bifurcaciones y dimensión,
  pp.~53--85; algoritmo de Lyapunov, pp.~211--216.
  \item Ott \cite{Ott2002Chaos}, pp.~168--210 y 304--344, y
  T{\'e}l--Gruiz \cite{TelGruiz2006Chaotic}, pp.~191--225, para sillas,
  escapes, fronteras y crisis.
  \item Caso primario de transitorio oculto
  \cite{Danca2016HiddenTransient}; exponentes fraccionarios y su algoritmo
  \cite{DancaKuznetsov2018Lyapunov}; memoria e inicialización
  \cite{Sabatier2021InfiniteMemory}.
\end{itemize}
\end{bloqueestudio}
\begin{bloqueestudio}{Ejercicios}
\textbf{Básicos.} Derive una cota de Gronwall para el defecto de una EDO y la
fórmula heurística del horizonte de predictibilidad. Calcule FTLE en tres
ventanas y muestre su distribución, no sólo la media.

\textbf{Investigación.} Fije una región de escape, estime la supervivencia
de un conjunto de semillas y compare una ley exponencial con una cola no
exponencial. Trate las trayectorias que no escaparon como datos censurados.

\textbf{Integrador.} Repita el diseño para Caputo en el espacio de historias:
defina las historias muestreadas, conserve el terminal inferior y separe
continuación de reinicio. Indique qué parte vive sólo en la proyección.
\end{bloqueestudio}
\fi}

\newcommand{\HAFOCardCases}{\ifHAFOPedagogy
\begin{bloqueestudio}{Lecturas}
Artículos primarios de los casos: Chua
\cite{LeonovKuznetsovVagaitsev2011,KuznetsovEtAl2017,
KuznetsovEtAl2023ChuaExperiments}; sistemas
fraccionarios \cite{Danca2017,Danca2017HiddenFO,Wu2023ArctanChua}; candidatos
no Chua \cite{munoz2018,wang2020}. Las revisiones
\cite{DudkowskiEtAl2016,GuanXie2025} se usan después para comparar métodos.
\end{bloqueestudio}
\begin{bloqueestudio}{Ejercicios}
\textbf{Básicos.} Recalcule equilibrios, jacobiano y clasificación local de un
caso entero y uno fraccionario.

\textbf{Investigación.} Para cada caso complete: campo; espacio de fases;
equilibrios; estabilidad; semilla; integración; destino; caos; cuenca;
conclusión autorizada. Compare en paralelo un caso entero y uno fraccionario:
qué objetos permanecen en el espacio de estados y cuáles pasan al espacio de
historias.
\end{bloqueestudio}
\fi}

\newcommand{\HAFOCardCandidates}{\ifHAFOPedagogy
\begin{bloqueestudio}{Lecturas}
\emph{Systems with Hidden Attractors}
\cite{PhamVolosKapitaniak2017Systems}: pp.~21--63 y 79--102; catálogo
\cite{PhamEtAl2018HiddenAttractors}; capítulo fraccionario, pp.~199--238,
en \cite{WangKuznetsovChen2021Multistability}.
\end{bloqueestudio}
\begin{bloqueestudio}{Ejercicios}
\textbf{Básicos.} Seleccione un candidato y recupere su artículo original.
Registre ecuaciones, parámetros, orden, inicialización y equilibrios.

\textbf{Integrador.} Calcule jacobiano y estabilidad local; use el criterio
entero o de Matignon según corresponda. Después diseñe una semilla y una prueba
de cuenca que separen atracción, caos y ocultedad.
\end{bloqueestudio}
\fi}

\newcommand{\HAFOCardFuture}{\ifHAFOPedagogy
\begin{bloqueestudio}{Lecturas complementarias}
\begin{itemize}[nosep]
  \item \emph{Topología y geometría diferencial}: conexiones, pp.~135--142;
  geodésicas, pp.~143--149; torsión y curvatura, pp.~151--165; métrica,
  pp.~165--186.
  \item \emph{Teoría geométrica de ecuaciones diferenciales II}:
  formas normales y bifurcaciones, pp.~185--244.
  \item Guckenheimer--Holmes \cite{GuckenheimerHolmes1983}: PDF aprox.~6--39,
  65--119 y 120--184.
  \item Jost \cite{Jost2005Dynamical}: conjuntos invariantes aislados,
  pares índice y Conley, pp.~70--97; entropía y dinámica simbólica,
  pp.~99--110. Esta lectura introduce herramientas avanzadas y no sustituye
  una prueba completa.
  \item \texttt{Barnsley - SuperFractals.pdf}: PDF~19--99, 201--395;
  Lorenz, PDF~9--30 y 116--207; Mandelbrot, PDF~33, 277 y desde 367.
  \item Diethelm \cite{Diethelm2010AnalysisFDE}: multi-término,
  pp.~167--186. Oldham--Spanier \cite{OldhamSpanier1974Fractional}:
  pp.~120--132, 148--164, 192 y 207--218.
\end{itemize}
\end{bloqueestudio}
\begin{bloqueestudio}{Ejercicios}
Elija un caso entero y uno Caputo. Formule para ambos un problema abierto sobre
variedades, retorno o frontera de cuenca; escriba el objeto geométrico, el
espacio de fases correspondiente, la prueba numérica y la hipótesis que aún
faltaría demostrar.
\end{bloqueestudio}
\fi}
